Un aspecto clave de diseño de las redes conmutadas es el relativo al encaminamiento. Esta función trata de encontrar rutas a través de la red entre pares de nodos finales comunicantes, de modo que la red se use de forma eficiente. Antes de comenzar con la monografia, revisemos los siguientes conceptos: Conmutación, Redes conmutadas y Encaminamiento.

\subsection{Conmutación}

Es el proceso de establecer una conexión física o lógica entre dos puntos de una red para que puedan intercambiar datos. Es como el trabajo que hace un operador de trenes al mover las vías para que el vagón tome el camino correcto hacia su destino. Sin la conmutación, necesitaríamos un cable directo entre cada dispositivo del mundo, lo cual sería imposible.

\subsection{Redes conmutadas}

Son redes en las que los datos, al viajar del emisor al receptor, pasan por una serie de nodos intermedios (conmutadores o switches). En estas redes, el camino no es fijo ni permanente; los recursos de la red se comparten y se asignan solo cuando es necesario transmitir. Existen dos tipos principales:

\begin{itemize}
	\item \textbf{Conmutación de circuitos}: Se reserva un camino exclusivo (como una llamada telefónica tradicional).
	\item \textbf{Conmutación de paquetes}: La información se divide en trozos pequeños que pueden viajar por rutas distintas (como Internet).
\end{itemize}

\subsection{Encaminamiento}

Es la inteligencia de la red. Mientras que la conmutación hace el trabajo de mover los datos, el encaminamiento es el proceso de seleccionar la mejor ruta a través de la red para que los datos lleguen a su destino final. Los dispositivos encargados de esto son los routers, que consultan tablas de rutas para decidir qué camino es el más rápido o el que está menos saturado en ese momento.

La función del encaminamiento es encortar rutas a través de la red entre pares de nodos que desean comunicarse. Esta implementado en la capa de red o capa 3 del modelo OSI. Muchas conexiones necesitaran una ruta a traves de mas de un conmutador, la cual debe tener las siguientes características:

\begin{itemize}
	\item \textbf{Eficiencia}: usar el mínimo de equipamiento posible
	\item \textbf{Flexibilidad}: debe proporcionar trafico razonable aun en condiciones adversas.
\end{itemize}